\documentclass{article}

\usepackage{amsmath}
\usepackage[inline]{enumitem}   % configure itemize and enumerate
\usepackage[normalem]{ulem}  % \uline{} for underlining terms
\usepackage{xurl}  % allow URLs to break anywhere for better wrapping
\usepackage{graphicx}
\usepackage{xcolor}
\usepackage{subcaption}
\usepackage{titlesec}  % for adjusting section spacing
\usepackage{minted}
% Table packages:
\usepackage{makecell}  % \Xhline
\usepackage{multirow}  % \multirow
\usepackage{diagbox}   % \diagbox
\usepackage{booktabs}
\usepackage{wrapfig}
% Use inconsolata font for monospaced text, which is not so wide.
\usepackage[varqu]{zi4}
\usepackage{newtxmath}

% hyperref makes hyperlinks in the resulting PDF.
% If your build breaks (sometimes temporarily if a hyperlink spans a page)
% please comment out the following usepackage line and replace
% \usepackage{mlsys2025} with \usepackage[nohyperref]{mlsys2025} above.
\usepackage{hyperref}

% Use the following line for the initial blind version submitted for review:
\usepackage{mlsys2025}
% If accepted, instead use the following line for the camera-ready submission:
% \usepackage[accepted]{mlsys2025}

% \raggedbottom  % Allow more flexible spacing at bottom of pages/columns
% \emergencystretch=2pt
% \clubpenalty=0
% \widowpenalty=0

\setlist[itemize]{leftmargin=1em, itemsep=0pt, topsep=2pt plus 2pt, parsep=2pt}
\setlist[enumerate]{leftmargin=1em, itemsep=0pt, topsep=2pt plus 2pt, parsep=2pt}
\setlength{\parskip}{0pt plus 0.5pt}
% \titlespacing*{\subsection}{0pt}{3pt plus 1pt minus 1pt}{2pt minus 1pt}
% \titlespacing*{\subsubsection}{0pt}{3pt plus 1pt minus 1pt}{2pt minus 1pt}
\setlength{\textfloatsep}{3pt plus 2pt minus 2pt}
\setlength{\floatsep}{3pt plus 2pt minus 2pt}
\setlength{\intextsep}{3pt}

\definecolor{mauve}{HTML}{9558B2}
\definecolor{rose}{HTML}{FF007F}

\newcommand{\XComment}[1]{}
\def\COMMENTSON{true}  % Set to false to disable all comments
\ifdefined\COMMENTSON
\newcommand{\yifan}[1]{\textcolor{rose}{Yifan: #1}}
\newcommand{\sasa}[1]{\textcolor{mauve}{#1}}
\newcommand{\yuchen}[1]{\textcolor{orange}{Yuchen: #1}}
\newcommand{\matei}[1]{\textcolor{purple}{Matei: #1}}
\newcommand{\egan}[1]{{\color{green} Egan: #1}}
\newcommand{\minsoo}[1]{{\color{red} Minsoo: #1}}
\newcommand{\needattn}[1]{{\color{blue} #1}}
\newcommand{\NA}[1]{{\color{blue} #1}}
\else
\newcommand{\yifan}[1]{}
\newcommand{\sasa}[1]{}
\newcommand{\yuchen}[1]{}
\newcommand{\matei}[1]{}
\newcommand{\egan}[1]{}
\newcommand{\minsoo}[1]{}
\newcommand{\needattn}[1]{#1}
\newcommand{\NA}[1]{#1}
\fi

\def\NAME{Nautilus}
\def\TITLE{Nautilus: An Auto-Scheduling Tensor Compiler for Efficient Tiled GPU Kernels}

\mlsystitlerunning{\TITLE{}}

\begin{document}

\twocolumn[
  \mlsystitle{\TITLE{}}
  \vspace*{-1em}
  \begin{mlsysauthorlist}
    \mlsysauthor{Yifan Zhao}{illinois}
    \mlsysauthor{Yuchen Yang}{illinois}
    \mlsysauthor{Matei Budiu}{illinois}
    \mlsysauthor{Sasa Misailovic}{illinois}
  \end{mlsysauthorlist}
  \mlsysaffiliation{illinois}{Siebel School of Computer Science, University of Illinois
  Urbana-Champaign, Urbana, IL, USA}
  \mlsyscorrespondingauthor{Yifan Zhao}{yifanz16@illinois.edu}
  \vspace*{1em}

  \begin{abstract}
    We present \NAME{}, a novel tensor compiler that moves toward fully automated
    \emph{math-to-kernel} optimization.
    \NAME{} compiles a high-level algebraic specification of tensor operators into
    efficient tiled GPU kernels.
    \NAME{}'s successive lowering design allows high-level optimizations, expression
    rewrites, and tile optimizations to be jointly applied in a single end-to-end system.
    \NAME{} presents a novel auto-scheduler that discovers sequences of high-level
    optimizations, while preserving the regular program structure
    needed by tile optimizers.
    \NAME{}'s auto-scheduler captures complex interactions and trade-offs in the high-level
    optimizations,
    including aggressive global transformations like advanced reduction fusion.

    \NAME{} is the first end-to-end tensor compiler capable of starting from a math-like
    description of attention
    and automatically discovering FlashAttention-3-like kernels,
    offloading the entire burden of optimization from the programmer to the compiler.
    Across five transformer-based models and 150 evaluation configurations on NVIDIA GH200 and RTX
    5090 GPUs, \NAME{} achieves up to 23\% higher throughput than state-of-the-art
    compilers on GH200
    and up to 42\% on RTX 5090, while matching or exceeding manually written cuDNN kernels
    on many \mbox{long-sequence configurations.}
  \end{abstract}
]
\typeout{parskip = \the\parskip}
\typeout{parindent = \the\parindent}
\typeout{baselineskip = \the\baselineskip}
\typeout{textfloatsep = \the\textfloatsep}
\typeout{floatsep = \the\floatsep}
\typeout{intextsep = \the\intextsep}
\typeout{abovedisplayskip = \the\abovedisplayskip}
\typeout{belowdisplayskip = \the\belowdisplayskip}

\printAffiliationsAndNotice{}

\section{Introduction}
{Modern attention-based machine learning (ML) models
are essential for language, vision, and other emerging tasks in artificial intelligence
\cite{gpt5,image-gen-paper,audio-gen-paper}. To efficiently deploy these models on GPUs,
system engineers need tools to bridge the performance gap
between a \emph{math-like model specification}, such as PyTorch or JAX code, and efficient
low-level code that maximally utilizes the underlying hardware.
Today, \NA{existing tools fall short of this goal on increasingly varied hardware:}
they are either too hard to use, requiring expertise in rapidly-evolving low-level
hardware details, or produce programs that are not sufficiently performant.

The most common approach today is writing manually optimized kernels in hardware vendor languages
like CUDA~\cite{cuda} and CUTLASS~\cite{cutlass}. They achieve high performance and have been
used in PyTorch~\cite{pytorch} and TensorFlow~\cite{tensorflow}.
However, programming in these low-level languages requires significant hardware and
algorithmic expertise and does not scale to rapidly evolving GPU architectures and ML models.
On the opposite side, schedule-based tensor compilers such as Halide~\cite{halide} and
TVM~\cite{tvm} optimize a high-level specification of the ML operator {in a math-like
tensor language.}
This design offloads the obligation of optimization from the programmer to the compiler,
which needs to search in a complex design space of \emph{schedules}
(sequence of transformations) to find an optimal kernel on the target GPU.
\NA{But this space mixes high-level decisions (e.g., fusion) with low-level ones
(e.g., thread mappings)}, and searching it is intractable; thus
today's high-level tensor compilers fail to match handwritten code in performance.

As a tentative trade-off between programmability and efficiency,
tile-based compilers and languages,
including Triton~\cite{triton}, Pallas~\cite{pallas}, and TileLang~\cite{tilelang},
have emerged as an abstraction layer over hardware vendor languages.
Here, a programmer manually expresses the ML operator as computation over tiles (fixed
size subtensors), and the tile compiler applies low-level SIMD
optimizations to map tile computation to the hardware.
With tile languages, a programmer still has to manually control many essential
high-level optimizations. This issue is especially evident for attention:
each version of FlashAttention
\cite{flash-attention,flash-attention-2,flash-decoding,flash-attention-3}
is an increasingly complex, manually written tile program
that must be significantly modified for every new GPU architecture.

\NA{We advocate for a \emph{math-to-kernel} compiler}:
the program\-mer writes operators in a math-like linear algebraic language
(e.g., attention as a sequence of matmuls and softmax),
and the compiler automatically finds high-performance kernels (e.g., FlashAttention).
\NA{Such a compiler is a viable solution that can yield ML
kernel efficiency competitive with manually-written code on modern GPU architectures.}
The central difficulty to materialize this goal is that we must bridge two types of optimizations
that live on different levels of abstraction:
(1)~operation-level optimizations that reason about compute expressions and data
dependencies of linear algebra operations,
and (2)~tile optimizations that map regular computation patterns to the underlying hardware.
\NA{Our key insight is that the two levels can be separated:
  tile compilers already perform (2) effectively and can serve as the compiler's
  low-level IR and backend, so schedules only need to capture (1),
  which drastically shrinks the space to search.
Automating (1) on top of tile compilers, however, raises two challenges.}
First, compiling many tensor operations into one or a few efficient kernels
requires new fusion techniques to handle complex data dependency.
In particular, \emph{advanced reduction fusion}, introduced by Neptune~\cite{neptune},
can fuse reduction loop nests, such as softmax and matmul in attention,
and it is a key optimization to allow compilers to automatically discover
FlashAttention from math-like attention.
\NA{Neptune already delegates tile optimizations to Triton,
but it relies on manually written schedules;}
it is still unclear how to apply such optimizations in a compiler without user intervention.
\NA{Second, no single intermediate representation (IR) suits all of these optimizations.
  Operation-level optimizations reason about both data dependencies (e.g., fusion)
  and compute expressions (e.g., expression rewriting):
  fusion needs loops and data dependencies,
  while expression rewriting needs tiles as values it can freely manipulate.
Tile compilers, in turn, need explicit tile memory accesses.}
% \sasa{We need to develop the argument more here. Why is tile promising and why is tile
% not a good language for programmers but an IR?}

\vspace{0.3em}
\noindent{\textbf{Our work: \NAME{}}. } \NAME{} is a novel tensor compiler for
attention-based ML models.
\NAME{} brings together operation-level optimizations
(advanced reduction fusion and expression rewriting)
and tile optimizations in one pipeline, to generate high-performance
GPU kernels on modern and emerging GPUs.
\NAME{} adopts the concept of schedules to regulate the legality of the sequence of transformations
and facilitate search-based exploration~\cite{tvm},
and automates the generation of schedules that incorporate advanced reduction fusion
via a novel auto-scheduler designed specifically to target tile compilers.
The auto-scheduler addresses three challenges that have kept such optimizations manual.
\NA{(1) \emph{Non-local fusion:} advanced fusion fuses two reductions together with all
  element-wise operations between them,
  so a pass that transforms one statement at a time cannot decide it locally.
  The auto-scheduler looks ahead along data dependencies to find these opportunities.
  (2) \emph{Tile-readiness:} operation-level transformations can leave the program in a
  form tile compilers handle poorly, e.g., advanced fusion can introduce scan dependencies.
  The auto-scheduler regularizes the program to keep it ``tile-ready'',
  i.e., easy to optimize by tile compilers.
  (3) \emph{Input-dependent trade-offs:} the best schedule to apply and tile compiler to use
  depend on the operator, input shape, and hardware.
The auto-scheduler generates promising schedules and lets the auto-tuner select among them.}

\NA{\NAME{} then successively lowers the scheduled program through three IRs,
  each at the abstraction level of one category of optimizations.
  Fusion and other scheduling happen in a scalar IR.
  Expression rewriting happens in a novel virtual-register tile IR (VR-tile IR),
  which expresses tiles as values rather than memory accesses.
  Tile optimizations happen in a memory-access tile IR (MA-tile IR),
  which \NAME{} translates to multiple tile compilers (Triton, Tawa~\cite{tawa}, and TileLang)
and selects the best-performing one.}

\NAME{} is the first end-to-end compiler to perform fully automated math-to-kernel
optimization on attention: \NA{starting from a simple linear algebraic representation of attention,
it automatically discovers variants of FlashAttention-3 and delivers FlashAttention-3 performance.}
%\sasa{be careful re Flashlight}
%\sasa{This is the first general end-to-end compiler system that can do this task. (we
% may need to massage it!)}

We evaluate \NAME{} on five \matei{needs updating}
% \sasa{[Is it five?]}
transformer-based models, with {150} total evaluation configurations, on two
modern GPU architectures: NVIDIA GH200 (Hopper) and NVIDIA RTX 5090 (Blackwell). We used
seven compiler backends and five hand-optimized libraries as baselines. \NAME{}-generated
attention kernels have up to 23\% higher throughput over the
state-of-the-art compiler systems on NVIDIA GH200 and up to 42\% on NVIDIA RTX 5090.
Moreover, our results are comparable to and for many configurations (with longer sequences)
faster than manually written cuDNN kernels.

\vspace{0.3em}
\noindent\textbf{Contributions:} The paper makes several contributions:
\begin{itemize}
  \item \NAME{}, a tensor compiler for attention-based ML models that fully automates the
    generation of high-per\-for\-mance GPU kernels from a program description (\S\ref{sec:design}).
  \item An auto-scheduler composing advanced reduction fusion and other high-level
    optimizations in \NAME{}'s scheduling language, while keeping programs tile-ready
    (\S\ref{sec:auto-sched}).
  \item A three-IR lowering pipeline: a scalar IR, a novel VR-tile IR for expression
    rewriting, and an MA-tile IR that enables multiple tile compilers as backends
    (\S\ref{sec:opt-pipeline}).
  \item An evaluation on multiple state-of-the-art attention models and GPU architectures,
    showing higher throughput than existing compilers and performance on par with
    hand-optimized kernels (\S\ref{sec:eval}).
\end{itemize}
}

\section{\NAME{} System Design}
\label{sec:design}
{Figure \ref{fig:design} shows the architecture of \NAME{}.
% \yifan{Tried to clarify what is taken from existing work vs.~homemade. See if that comes across.}
% \matei{Idea: maybe we can add to the caption of Fig. 1 something like ``Our contributions
%   are shown in color'' or something like that. Or if that's not accurate, outline the boxes
% that we did in some color and make the arrows we did (translations) that color too.}
\NA{The input to \NAME{} is an operator written in TVM's tensor expression (TE) language,
  either directly or converted from PyTorch or TensorFlow layers by TVM Relax~\cite{relax}.
  First, \NAME{}'s auto-scheduler (\S\ref{sec:auto-sched}) finds high-performance schedules
  for the program on the target GPU without user intervention,
  while keeping the program \emph{tile-ready}, i.e., in the regular form tile compilers expect,
  with perfectly nested tile-sized inner loops and no scan dependencies.
  Schedules are written in a scheduling language over the program's \emph{block graph}
  (\S\ref{sec:block-graph}).
  Then, \NAME{} successively lowers the scheduled program from its scalar IR
  through two tile IRs (\S\ref{sec:opt-pipeline}):
  a virtual register tile IR (VR-tile IR) for expression rewriting,
  and a memory access tile IR (MA-tile IR), which \NAME{} translates into the language of
  a tile compiler (Triton, Tawa, or TileLang) that performs tile-level SIMD optimizations
and generates the kernel.}
Besides TE, other tensor expression languages, e.g., JAX's jaxpr,
\mbox{could also serve as \NAME{}'s input.}

\begin{figure}[t]
  \centering
  \includegraphics[width=0.8\columnwidth, trim=0mm 61mm 2mm 0mm, clip,
  page=1]{figures/one-col-figs.pdf}
  \vspace{-.05in}
  \caption{System overview of \NAME{}. Novel components from our contributions are shown
  in colored boxes.}
  % \sasa{We should either remove Tiled optimizers box or call it 'Codegen to Tile Optimizers'}}
  \label{fig:design}
\end{figure}

\begin{figure}[t]
  \centering
  \includegraphics[width=0.9\columnwidth, trim=0mm 5mm 25mm 0mm, clip,
  page=2]{figures/one-col-figs.pdf}
  \vspace{-.1in}
  \caption{An example of scheduling a block graph.
    A scalar IR program and its corresponding block graph are shown at the top,
    followed by examples of using scheduling primitives to transform the block graph. In
  this graph, ovals represent \emph{compute statements} and rectangles represent \emph{loops}. }
  \label{fig:sched-example}
\end{figure}

\textbf{Block Graph.}
\label{sec:block-graph}
The block graph is the program representation that \NAME{}'s scheduling manipulates.
It combines the AST of a scalar IR program with its data dependency graph,
\NA{and models the program only down to the granularity of a GPU thread block:
  which loops and statements each block executes, how data flows between statements, and
  where tensors are placed. Lower-level details, such as how a block's tile
  computation maps to warps and threads, are absent from the block graph and left to tile compilers.
  % (unlike TensorIR schedules~\cite{tensorir}, which also bind threads and tensorize).
This keeps schedules high-level and the auto-scheduler's search space small.}
The top part of Figure \ref{fig:sched-example} shows an example scalar IR program
(extracted from the attention operator) and its block graph.
Each node corresponds to a \emph{compute statement} in the IR,
and directed edges indicate data dependencies.
The \emph{producers} of a compute statement $n$ are the statements with an edge to $n$,
and $n$ is their \emph{consumer};
\emph{predecessors} and \emph{successors} are \NA{the transitive closures of these relations},
and a \emph{path} from $n_1$ to $n_2$ is a chain of statements, each a producer of the next,
from $n_1$ to $n_2$.
Paths are central to advanced fusion,
which fuses two reductions together with all statements on the path between them
(\S\ref{sec:adv-fusion}).
Boxes around compute statements represent \emph{loops}.
% \matei{term unclear. Intended: spatial/reduction/scan.}
The block graph tracks iteration-carried data dependency of each loop,
determined by the data access pattern of compute statements.
For example, the \texttt{k} loop around the \texttt{mm} node is a reduction loop,
while the loops \texttt{i} and \texttt{j} are map loops.

\NA{\textbf{Scheduling Language.}}
\NAME{}'s scheduling language consists of primitives that transform the block graph
(Table~\ref{tab:primitives} in the appendix).
They cover loop tiling and binding loops to GPU thread blocks,
data placement and caching,
classic loop fusion and advanced reduction fusion (rolling update and split-k update),
and transformations that regularize the program for tile compilers,
such as decomposing reductions and splitting scan buffers.
Many primitives come from TVM and Neptune;
\NAME{} adds primitives for padding axes, folding accumulators, and reordering statements.
The bottom part of Figure~\ref{fig:sched-example} shows examples of applying primitives
to the block graph.
}

% \vspace{-.05in}
% \section{Background}
% \label{sec:background}
% {\input{03-background.tex}}

\section{\NAME{} Auto-Scheduler}
\label{sec:auto-sched}
{The \NAME{} auto-scheduler finds high-performance schedules for an input program without
user intervention. \NA{Three challenges shape its design.
  First, advanced fusion is non-local: it fuses two reductions together with all statements
  on the path between them, so it cannot be decided one statement at a time;
  the auto-scheduler looks ahead along data dependencies to find these opportunities
  (\S\ref{sec:adv-fusion}).
  Second, fusion can leave the program in forms that tile compilers optimize poorly,
  such as scan dependencies;
  the auto-scheduler regularizes the program to keep it tile-ready (\S\ref{sec:regularize}).
  Third, the best schedule and tile compilers depend on the operator, input shape, and GPU;
  the auto-scheduler sketches a set of promising schedules,
  which we call \textit{schedule seeds},
  and leaves their decision parameters to the auto-tuner (\S\ref{sec:auto-tuner}),
while heuristics prune the seeds aggressively to keep auto-tuning \mbox{within minutes.}}

\subsection{Auto-scheduling Algorithm}\label{sec:autosched-alg}
The auto-scheduler builds up schedule seeds by applying a sequence of \emph{scheduling passes}
(Table~\ref{tab:passes}), each applying one class of optimizations such as tiling or fusion.
% The input to the auto-scheduler is a program in which all compute statements
% are perfectly nested in their own loops.
It visits compute statements in a pre-order traversal, producers to consumers,
while maintaining a list of candidate schedules.
At each statement, it applies the passes in order to every candidate;
a pass may replace a candidate with several if there are multiple promising ways
to transform the statement, keeping only those its heuristics favor
(e.g., fusing into the deepest feasible loop).
After the last statement, the final two passes run on each finished candidate,
and the candidates become the schedule seeds,
whose design parameters (e.g., block tiling sizes) the auto-tuner explores.

\begin{table*}[t]
  \centering
  \small
  \caption{Compile passes in the Nautilus auto-scheduler, in the order they are applied.
    The last two run on
  each finished candidate.}
  \label{tab:passes}
  \begin{tabular}{lll}
    \hline
    Scheduling Pass & Transformations & Description \\
    & Used & \\
    \hline
    PadRows            & 1                & Adds a candidate with short query rows
    (decode, MTP) padded to 16 for tensor cores \\
    InlineReshapes     & 1                & Inlines tensor reshaping \\
    AdvTilingFusion    & 4   & Uses advanced Neptune fusion to fuse compute statements \\
    CacheSetup         & 4  & Shares or rematerializes large intermediates per tile in
    shared memory \\
    DecomposeSplitScan & 3      & Binds grid loops; apply decompose reduction or split
    scan buffer \\
    & & to every compute statement which had advanced fusion \\
    BackendSelection   & (none)            & Configures compile parameters for different backends \\
    BlockOrder         & 1                & Makes the order of statements in each loop a
    tuning choice \\
    DataPlacement      & 4   & Places every tensor in global, shared memory or registers;
    caches index values \\
    Finalize           & 2          & Turns masks into loop ranges and folds accumulators
    before lowering to tile IR \\
    \hline
  \end{tabular}
\end{table*}

\subsection{Auto-Scheduler Transformations}
Our auto-scheduler applies four families of transformations:
tiling, fusion, cache management, and dataflow management.
These transformations capture the main steps needed to translate
and optimize a tensor-expression operator into a tile program ready for \NAME{}'s tile optimizers.
% A family of transformations may be implemented as multiple passes in the auto-scheduler
% to break dependencies between certain transformations. \matei{It might be unnecessary
% to talk about
% passes at all if we give an idea of the transformation types.}

\subsubsection{Tiling.}
The bi-level tiling transformation applies to a perfect loop nest,
splitting each loop into an inner loop and an outer loop,
and reordering them so that all outer loops are outside of all inner loops.
For example, a loop nest with two loops \texttt{i, j} turns into the loops
\texttt{i0, j0, i1, j1}, with the first two loops being the outer loops.
The other passes in the auto-scheduler transform only the outer loops;
the inner loops remain perfectly nested so they cleanly map to tiles in the tile IRs later.
Later, outer loops that are parallelizable (no data dependency across iterations)
are mapped to GPU blocks.
A size limit on inner loops is also enforced to prevent tiles from becoming
too large and not fitting into the cache - in this case, the respective large inner loop is itself
tiled into an inner and outer loop if possible.

% \yifan{Good to have but I'm struggling where to put this, maybe not here.}
% \matei{moved to dataflow}

\subsubsection{Operator Fusion.}\label{sec:adv-fusion}
Fusion reduces the number of kernels and improves data locality.
Each fusion primitive fuses a compute statement into a loop nest.
\textbf{Classic loop fusion} preserves data dependencies
and fuses a statement into the element-wise loop nest of a direct producer.
Fusing into a reduction loop nest requires \textbf{advanced fusion}~\cite{neptune},
which fuses reductions into an ``online'' form by intentionally breaking some dependencies,
and repairing the program with changed compute expressions, at the cost of extra computation.
Advanced fusion applies to a pair of reductions,
but also needs to capture all element-wise statements on the path between them.

\NAME{} has two advanced fusion primitives.
\textbf{Rolling update} fuses the reductions into a single kernel.
\textbf{Split-k update} computes a partial result for each iteration of the reduction loop
in parallel thread blocks, then combines the partials in a second kernel with the same repair;
this gains parallelism along the reduction axis when the other loops cannot fill the GPU,
as happens in decoding attention.
The auto-scheduler is capable of \textbf{hybrid strategies}, applying split-k update across chunks
of the reduction and rolling update within each chunk.
% which suits kernels with few query rows but a long reduction,
% such as multi-token prediction verification.
% The auto-scheduler chooses the strategy when tiling the loop nest.

For each program statement $s$, the fusion pass first greedily applies classic fusion
into the loop nests of its direct producers, at the deepest possible level.
% \yifan{I moved this over from ``additional considerations'' -- it can be beneficial, but
% what do we actually do?} \matei{not supported yet to fuse with multiple producers, we try
% either one as a candidate}
When a statement has multiple producers, the pass tries fusing with each one and produces
multiple candidate schedules.
If $s$ is a reduction, the pass also tries advanced fusion into each predecessor
reduction, which, when feasible, always fuses deeper than classic fusion.
If an element-wise $s$ cannot be fused, the pass looks ahead for a successor reduction
whose advanced fusion would capture $s$; if none exists, $s$ starts a new kernel.

Consider a loop nest $L$ with an element-wise loop $e$ around a reduction loop $r$.
$L$ has an element-wise consumer $c_1$, and $c_1$ has a reduction consumer $c_2$.
On its own, $c_1$ can only fuse into $e$ via classic fusion.
Advanced fusion of $c_2$ into $r$ would also capture $c_1$ and place it under $r$,
deeper than classic fusion can, but this opportunity is only visible when examining $c_2$;
this is what makes advanced fusion \emph{non-local}.
Since pre-order traversal visits $c_1$ first, the fusion pass greedily fuses $c_1$ into $e$.
This does not prevent the deeper fusion: when the pass later fuses $c_2$ into $r$,
it fuses $c_1$ again, moving it from under $e$ to under $r$.
Therefore, \NAME{}'s fusion pass handles non-local advanced fusion with a plain
pre-order traversal, without anticipating advanced fusion when making classic fusion decisions.

% Matei: moved earlier | and a ``rolling update'' technique that enables fusion between reductions.
% Rolling update is an advanced fusion technique from Neptune \cite{neptune}
% which is beyond the capabilities of classic loop fusion,
% fusing multiple reductions into an ``online'' kernel.
% While classic loop fusion preserves data dependency,
% rolling update intentionally breaks some dependency constraints
% and repairs the program by changing compute expressions.
% \matei{the greedy heuristic part might make more sense to include in section 4.4
% instead of here if we are discussing advanced fusion there}
% and in the case of RollingUpdate, we prefer to capture more loop nests.
% In the case where there isn't a Pareto-dominant choice (such as when one fusion choice
% is deeper but
% the other captures more loop nests), both strategies are considered as produced schedules.
% The fused loop nests do not have to be
% adjacent, in which case any element-wise loop nests along the data dependency path
% between the target loop nests
% will be `captured' and fused as well
% This pass operates in a pre-order traversal of the schedule. The program's compute statements are
% traversed from producer to consumer, and existing candidate schedules are evaluated for
% fusion opportunities
% at the compute statement - when fusion is possible with one of that compute statement's loops,
% new candidate schedules are added, replacing the original unfused candidate.

\subsubsection{Cache Management.}
Cache management places tensors on an inner (more local) memory level,
i.e., shared memory or registers rather than global HBM, to reduce slow global memory traffic.
All tensors start in global memory. The auto-scheduler performs two types of data
localization: (1) it caches a copy of a tensor to shared memory if \NA{a reduce loop
nest} exhibits temporal reuse of the tensor, and (2) it places the tensor itself
in an inner memory level if all of its accesses are local to a thread block.
Tensors that are not explicitly cached are placed by how they are used:
in global memory if another kernel reads them (such as split-k partials),
\NA{in shared memory for matmul operands whose layout is incompatible with
  tensor core matmul, and in registers otherwise.
  A tensor core matmul has strict requirement for the layout of its register operands,
  which can cause tile compilers to fail on programs.
  \NAME{} anticipates layout failures using heuristic rules and places tensors on shared memory
when a relayout is needed.}
% \NA{ A matmul's second
%   operand always goes to shared memory: every warp reads all of its columns, so keeping it
%   in registers would require redistributing it across threads before each matmul. The first
%   operand may be the result of an earlier matmul, held in registers and spread across
%   threads in that matmul's output layout. It can stay there only if this layout matches the
%   operand layout the second matmul expects, which holds under the following conditions: (1)
%   The second matmul reads the tile along its reduction axis. (2) The tile has at least 16
%   rows per warp (128 rows for the 8 warps we can sample), so both matmuls split rows across
%   warps the same way. With fewer rows, threads hold elements the consumer expects in other
%   threads, and moving them would need a cross-thread exchange that the tile optimizer does
%   not insert. (3) Finally, the operand is not FP8. FP8 matrix instructions pack four
%   consecutive reduction-axis elements per register, while the FP32 accumulator gives each
% thread pairs of adjacent columns, so the layouts will never match.} \matei{this may be
% too detailed...}
% Later in the tile optimizer, tensors whose register layout does not match what their
% consumer expects are rejected, so cache management places these tensors in shared memory
% such that the tile optimizer can reload them with a relayout.
% ; this applies to a matmul's second operand, and an 8-bit operand computed in registers,
% which inherits its FP32 accumulator's layout while FP8 tensor cores expect a different one.

Inner memory levels such as shared memory are small,
so the auto-scheduler analyzes tensor sizes to identify
large buffers with multiple readers that would otherwise occupy scarce local memory.
If such a buffer's producer and all its readers share a loop, and the producer's
tile at that loop fits in shared memory, the auto-scheduler keeps one copy of
that tile for all readers. Otherwise, it rematerializes the buffer,
recomputing the buffer near each use to reduce memory pressure and improve occupancy.
Integer tensors used for indexing, such as sparse attention's top-k indices,
are placed by how they are accessed.
\NAME{} tries to eliminate integer intermediates used only for indexing (in a gather
operation) by inlining their computation into their consumers.

% Each thread then loads the indices it needs directly from global memory, instead of being
% restricted to those held in its own registers, so the gather can be spread across all warps.
% \NA{When the auto-scheduler folds accumulators,} an integer input used both in gathers
% and as a value in computation gets a shared memory copy for the value consumer only,
% staged once per tile and prefetched with the gathers, while the gathers still read it
% directly from global memory.

\subsubsection{Dataflow Management.}\label{sec:regularize}
Our auto-scheduler applies several transformations to align the data dependency patterns
of the program with what tile compilers expect,
as tile compilers often optimize only a small set of tile computation patterns well.
One pattern they handle poorly is the ``scan'' dependency
(e.g., \texttt{a[i, j] = f(a[i - 1, j])}),
which can exist in the input program itself
or arise from previous transformations, such as rolling update;
existing tile compilers often fail to optimize scan dependencies
or produce significantly suboptimal kernels.
The auto-scheduler removes them with a split-scan-buffer transformation,
which multi-buffers the scan's output and updates the buffers once per loop iteration.
\NAME{} uses range analysis to convert expressions of boolean masks
into tightened loop ranges that skip fully masked tiles,
which improves the performance of masked attention, such as causal attention.
\NA{Because matmuls with fewer than 16 rows cannot use tensor cores,
  \NAME{} also adds candidates that pad such row axes to 16 (e.g., query rows in decoding),
zero-filling the computation and masking the load and stores.}

% To maintain correctness, if a value computed inside a rolling update loop is read
% from outside of that loop, the candidate is rejected, as the value would not be properly rescaled.
% \yifan{Sentence moved from Auto-Tuner.}

\input{figures/three-ir-ex.tex}

\subsection{Auto-Tuner}\label{sec:auto-tuner}
\NAME{} uses MetaSchedule's auto-tuner~\cite{metaschedule} to search the decision parameters
of the schedule seeds.
Every tunable decision is a sampled variable in the schedule trace,
which MetaSchedule's evolutionary search mutates and child schedules inherit:
the search generates schedules by varying these decisions, predicts their performance
with a learned cost model, and measures the projected-best schedules on the target GPU.
The tuning budget is the number of GPU measurements.
% \matei{The search space is pruned by dropping any seeds which fail to build.
%   Many fusion structures can't be built with any tile sizes - these would otherwise dominate
%   the initial tuning population.
%   \textit{note: is this something to include? might make people question why we are
% producing invalid seeds?}}

\NA{\textbf{Tunable decisions.}
  A full list of tunable decisions is given in Appendix Table~\ref{tab:decisions}.
  Structural decisions become separate seeds during auto-scheduling (\S\ref{sec:autosched-alg}),
  such as fusion strategy, whether row padding is applied, and the tile compiler to use.
  The BackendSelection pass splits each seed into one per tile compiler (Triton,
  Tawa, or TileLang).
  Within each seed, there are integer and categorical knobs to be tuned.
  Some apply directly to the schedule, such as tile sizes and the KV split count.
  The order of certain expensive statements (matmuls, reductions, copies) is also tunable.
  Other knobs become annotations on the IR program and take effect in later stages:
  switchable expression rewrites (such as accumulator folding),
  and tile optimizer parameters (e.g., TMA copy flags).
  % BackendSelection records these as annotations that the implementing stage reads.
  This design allows the auto-tuner to explore all these decisions from multiple levels,
jointly making decisions on tiling, fusion, expression rewrite, and tile optimization.}
% \matei{determine if table will be included}}

\section{\NAME{}'s Successive Lowering Pipeline}
\label{sec:opt-pipeline}

After scheduling, \NAME{} lowers the program from its scalar IR through two tile IRs,
because no single IR suits both remaining categories of
optimizations: expression rewriting needs tiles as values that it can freely combine and move
across statements and loop iterations, while tile compilers need explicit tile loads and stores
to perform SIMD optimizations.
The virtual register tile IR (VR-tile IR) serves the former and is new in \NAME{};
the memory access tile IR (MA-tile IR) serves the latter and is purposefully similar to
the tile IRs of existing tile compilers, simplifying interfacing with them.

\textbf{Tile IRs.}
Figure \ref{fig:3ir-example} shows an example program in both tile IRs,
and scalar IR for reference.
The MA-tile IR exposes tile data movement as explicit tensor accesses,
which is typical in tile IRs of existing tile compilers.
\NAME{}'s VR-tile IR expresses tiles as local variables (``virtual registers'')
rather than as memory accesses, similar in spirit to memory-to-register promotion
(``mem2reg'') in compilers such as LLVM.
This design reduces data movement instructions that would otherwise obstruct expression
manipulation, and enable powerful rewrites.
For example, \texttt{vmm}, \texttt{vmax\_local}, and \texttt{vmax} in Figure \ref{fig:vr-ir-example}
are each bound to a tile expression, so expression rewrites view them as one
expression when needed and freely move terms across them.

Virtual registers cover values within a loop body, but a value that evolves across loop
iterations, such as a running maximum, would still be expressed as repeated update-writes
to a tensor, hiding it from expression rewrites.
The VR-tile IR therefore includes a for-loop expression, which takes loop-carried
parameters and returns values: in Figure \ref{fig:vr-ir-example}, the loop expression
\texttt{for j0 in loop(M0)}
takes parameter \texttt{vmax\_1} and returns \texttt{vmax}.
Loops with loop-carried values also exist in other IRs, e.g., MLIR's \texttt{scf.for}~\cite{mlir};
the VR-tile IR uses them over tiles, so that tile value evolution over loop iterations
is expressed natively and tensor expression rewriting can operate across iterations.

\textbf{Lowering to VR-tile IR.}
\NAME{} lowers the scheduled scalar IR into the VR-tile IR,
relying on the tile-ready form that the auto-scheduler maintains (\S\ref{sec:auto-sched}).
Each perfectly nested inner loop nest becomes a tile expression bound to a virtual register,
such as the \texttt{dot} for the matmul in Figure~\ref{fig:vr-ir-example};
an outer loop that updates a tensor across its iterations, such as the running maximum
\texttt{xmax}, becomes a loop expression whose loop-carried parameter (\texttt{vmax\_1})
holds that tensor's tile.

\textbf{Expression Rewriting.}
\NAME{}'s expression rewriter applies a set of rewrite rules recursively on the VR-tile IR.
Table \ref{tab:rewrite-rules} lists representative rules, including rules over loop expressions
and rules informed by GPU hardware (e.g., rewriting $e^x$ using $2^x$).
These rules matter more after advanced fusion,
whose repair computation (\S\ref{sec:adv-fusion}) introduces rescaling in every loop iteration.
For example, the fourth rule in Table \ref{tab:rewrite-rules}
cancels out division and multiplication over loop iterations, removing $2n - 1$
operations from the program,
and the second rule computes the online-softmax rescale once per row instead of per element.
Expression rewriting reduces the latency of \NAME{}'s kernels in our ablation study
(Table~\ref{tab:ablation}).

\begin{table}[ht]
  \centering
  \scriptsize
  \setlength{\tabcolsep}{3pt}
  \renewcommand{\arraystretch}{1.1}
  \caption{Representative rules in \NAME{}'s expression rewriter. Standard algebraic
  identities are used but not listed.}
  \tt
  \label{tab:rewrite-rules}
  \begin{tabular}{l}
    \hline
    \textbf{Pattern} $\Rightarrow$ \textbf{Replacement} \\
    \hline
    $ \exp(x) \Rightarrow \mathrm{exp2}(\mathrm{log2}(e) \cdot x) $ \\[2pt]
    $\mathrm{exp2}(\text{bcast}(a) - \text{bcast}(b))$ $\Rightarrow$
    $\text{bcast}(\mathrm{exp2}(a - b))$ \\[2pt]
    $ x = \text{for}(...)\ \{ \text{yield } c \cdot x \} $ \\
    \quad $\Rightarrow$ $ x' = \text{for}(...)\ \{ \text{yield } x \};\ x = c \cdot x' $ \\[2pt]
    $ x = \text{for}(i: n)\ \{ \text{yield } x / y_i \cdot y_{i+1} \} $ \\
    \quad $\Rightarrow$ $ x' = y_n \cdot \text{for}(i: n)\ \{ \text{yield } x \} $ \\[2pt]
    $\text{select}(c, \text{dot}(A,B) \cdot s, k)$ \\
    \quad $\Rightarrow$ $\text{dot}(A, B, \text{init} = \text{select}(c, 0, k/s))\cdot s$ \\
    \hline
  \end{tabular}
\end{table}

\textbf{Tile Optimizations.}
After expression rewriting, \NAME{} translates the program into the MA-tile IR
by desugaring loop expressions and tile-valued virtual registers
into plain for-loops and explicit tile loads and stores.
The MA-tile IR builds on the tile IR of Neptune, which Neptune lowers to Triton.
We design the MA-tile IR as a superset of multiple existing tile languages,
which enables \NAME{} to use multiple tile compilers as its tile optimizers
(Triton, Tawa~\cite{tawa}, and TileLang~\cite{tilelang}), as chosen by the auto-tuner
(\S\ref{sec:auto-tuner}).
}

\vspace{-.08in}
\section{Implementation}
\label{sec:impl}
{\NAME{} is implemented on TVM's TensorIR infrastructure~\cite{tvm}:
its scheduling language builds on TensorIR's, and all three IRs use TensorIR code
infrastructure as a foundation.
The scalar IR and the MA-tile IR are adapted from Neptune~\cite{neptune},
while the VR-tile IR is new.
Each scheduling pass is implemented as a TVM MetaSchedule rule, driven by a custom
design-space generator that implements the pre-order traversal of \S\ref{sec:auto-sched}.
A runtime manages kernel compilation, caching, and invocation for all tile compilers
through a unified interface.

%\sasa{We should use stronger wording here so it doesn't sound like we are just lego bricks guys.}
%\sasa{Say what we implemented in autoscheduler :) }
%is implemented on top of TVM's MetaSchedule framework, but with entirely custom components. We 
%
%The \NAME{} auto-scheduler has an auto-tuning stage that leverages MetaSchedule's
% (Auto-tuner search algorithm and budget are described in Sec. 3.4.)
}

\section{Experimental Methodology}
\label{sec:method}
{\textbf{Hardware Platforms.}
We evaluate \NAME{} on an NVIDIA GH200 Grace Hopper Superchip (H100 GPU, Hopper)
and an NVIDIA RTX 5090 GPU (Blackwell).

\textbf{Workloads.}
We evaluate the attention operators of five transformer models
(Appendix Table~\ref{tab:operators}):
global attention from ViT 1.2B, causal attention (MHA) from Llama 2 7B, Qwen2 7B, and Qwen3 8B,
and grouped-query attention (GQA) from GLM-4 9B.
In end-to-end experiments, we run each model in prefill mode
on batch sizes 1 and 8, sequence lengths $2^{10}, \dots, 2^{15}$,
and FP16 and FP8-E4M3 precisions;
in operator experiments, we evaluate the three attention operators individually
on the same inputs.
We refer to the tuple of hardware platform, operator, input shape, and input precision as
a \textbf{setup}.
For each setup, we run \NAME{}'s auto-tuner for 256 empirical measurements.

\textbf{Baselines.}
We compare \NAME{} to 12 baselines (names we use are underlined):
seven compilers, namely \uline{Tawa}~\cite{tawa} and \uline{TileLang}~\cite{tilelang}
with their FlashAttention-3-style attention kernels,
two Triton implementations (\uline{OpenAI Triton}~\cite{openai-fused-attn}
and \uline{Tri-Dao Triton}~\cite{flash-attn-repo}), \uline{Helion},
FlexAttention (\uline{FlexAttn}) in PyTorch, and \uline{TVM};
and five hand-optimized libraries, namely \uline{cuDNN},
Dao-AILab's CUTLASS FlashAttention (\uline{Tri-Dao})~\cite{flash-attn-repo},
\uline{FlashInfer}~\cite{flashinfer}, \uline{ThunderKittens}~\cite{thunderkittens},
and \uline{PyTorch}'s built-in attention~\cite{pytorch}.
Appendix Table~\ref{tab:baselines} lists versions and configurations.
We do not evaluate FlashAttention-4~\cite{zadouri2026flashattention4algorithmkernelpipelining}
because it uses approximation techniques outside the scope of exact attention.

\textbf{Profiling and Performance Reporting.}
We report throughput in TFLOPs per second, dividing the theoretical attention FLOPs
of each configuration~\cite{Megatron-LM} by the measured latency,
which we measure with CUDA Events in end-to-end experiments
and with Nsight Systems in operator experiments (Appendix~\ref{sec:exp-details}).
We warm up once, report the mean latency of 10 measurements,
and average across setups with geometric means.

% \yuchen{GH200 has H100 as GPU} \yifan{thanks for the correction!}
%\sasa{Motivate that these are sota hw architectures with modern features...}
% \NA{Sparse} & \NA{Deepseek Sparse Attention} & \NA{DeepSeek-V3.2} \\
%\sasa{Can we motivate why these networks are representative?}
% Each operator is evaluated in both prefill and decoding modes:
% in prefill (PF), the query sequence length $s_q$ equals the key/value sequence length $s_{kv}$,
% and in decoding (DC), $s_q = 1$.
% Variations on masking (Global and Causal) collapse into the same operator under decoding
% mode, so we select one.
% Each operator is evaluated in prefill (PF) mode, where the query sequence length $s_q$
% equals the key/value sequence length $s_{kv}$.
%\yifan{Reminder: ThunderKittens kernel is not ``official'' but provided by one of the maintainers.}
% \sasa{Doublechk; cite the method if we have from somewhere}
%We also report \textbf{relative performance} of a baseline relative to \NAME{},
%which is $t_0 / t_i$ for the $i$-th baseline.
}

\section{Evaluation}
\label{sec:eval}
{% \yifan{Rerun evaluation with updated baselines. Update figures 5-8 and Tables 5 and 6.
% Table 7 should be good.}

% \yifan{Add results on Sparse Attn (with Deepseek 3.2); some multi-modal LLMs.}

\noindent\textbf{RQ1:}
Can \NAME{} accelerate end-to-end transformer-based model inference and outperform existing
tensor compilers and attention-specialized inference frameworks?
%Does \NAME{}'s performance benefit generalize to a variety of operators and input shapes?

\begin{figure*}[t!]
  \centering
  \includegraphics[width=0.9\textwidth]{figures/throughput_plot.pdf}
  \vspace{-1em}
  \caption{Throughput (TFLOPs/s) on GH200 across different models.
    We evaluate \NAME{} against baselines across five foundation models (GLM,
    Llama, Qwen2, Qwen3, ViT).
    The benchmarks cover FP16 and FP8 precisions with Batch Size 1 and 8, for sequence
  lengths ranging from 1k to 32k. }
  %\matei{This might be really annoying to do, but if it isn't too big a deal, I would
  % make it so that the bars are colored a bit differently (maybe with patterns) so they
  % can be easily distinguished in grayscale. I might also make the Nautilus color a bit
  % more distinct in some way to make it pop out more.}\yuchen{Of course we can change
  % the colors, any suggestions?}}
  \label{fig:gh200-throughput}
\end{figure*}

\begin{figure*}[t!]
  \centering
  \includegraphics[width=0.95\textwidth]{figures/throughput_plot_5090.pdf}
  \vspace{-1em}
  \caption{Throughput (TFLOPs/s) on RTX 5090 across different models.
    We evaluate \NAME{} against baselines across five foundation models (GLM,
    Llama, Qwen2, Qwen3, ViT).
    The benchmarks cover sequence lengths ranging from 1k to 32k.
    % \sasa{How do we distinguish between batch 1 and 4?}
  }
  \label{fig:5090-throughput}
\end{figure*}

\noindent\textbf{RQ2:}
Can \NAME{} generate kernels that surpass the performance of kernels from existing tensor compilers?

\noindent\textbf{RQ3:}
Can \NAME{} generate kernels that have comparable performance to manually optimized
kernels in inference frameworks, which represent today's state of the art?

\noindent\textbf{RQ4:}
Ablation studies: How much do components of \NAME{} contribute to the performance of
\NAME{}'s attentions?

\subsection{End-to-End Model Inference}
\label{sec:eval-e2e}

\begin{table}[b!]
  \centering
  \small
  \setlength{\tabcolsep}{2.5pt}
  \caption{Geomean of relative throughput of \NAME{} against baselines on GH200
  across sequence lengths and batch sizes.}
  \label{tab:rq1-summary}
  \def\dgbox{\diagbox[width=6em, height=2.5em]{Model}{Method}}
  \renewcommand{\arraystretch}{0.93}
  \begin{tabular}{l|cccccc}
    \toprule
    \dgbox & Flash-2 & Flex & SDPA & Tawa & Flex-FP8 & Tawa-FP8 \\
    \midrule
    GLM     & 1.17 & 1.03 & 1.33 & 1.11 & 1.07 & 1.16 \\
    LLAMA   & 1.23 & 1.06 & 1.26 & 1.18 & 1.08 & 1.17 \\
    QWEN2   & 1.20 & 1.05 & 1.20 & 1.12 & 1.07 & 1.13 \\
    QWEN3   & 1.21 & 1.05 & 1.21 & 1.12 & 1.06 & 1.13 \\
    VIT     & 1.28 & 1.54 & 1.15 & 1.28 & 1.59 & 1.43 \\
    \bottomrule
    \textbf{Average} & 1.22 & 1.13 & 1.23 & 1.16 & 1.16 & 1.20 \\
    \bottomrule
  \end{tabular}\vspace{-.03in}
\end{table}

\textbf{Overall Trends.}
% \sasa{Aren't we talking about throughput here?}
To answer RQ1, we evaluate \NAME{} on two hardware generations: GH200 (Hopper) and RTX
5090 (Blackwell), using five representative foundation models:
GLM~\cite{du2022glmgenerallanguagemodel}, Llama2~\cite{touvron2023llama2openfoundation},
Qwen2~\cite{yang2024qwen2technicalreport}, Qwen3~\cite{yang2025qwen3technicalreport}, and
ViT~\cite{dosovitskiy2021imageworth16x16words}.
Table~\ref{tab:rq1-summary} and Table~\ref{tab:rq1-summary-5090} present the averaged
relative throughput of \NAME{} relative to state-of-the-art baselines including both
vendor-optimized libraries (FlashAttention-2, SDPA) and compiler-based approaches
(FlexAttn, Tawa~\cite{tawa}), where each cell represents the geometric mean over varying
sequence lengths and batch sizes.

Across the five representative foundation models in Table~\ref{tab:rq1-summary}, \NAME{}
achieves consistent performance gains in both FP16 and FP8 regimes.
In the FP16 setting, \NAME{} achieves an average speedup of $1.22\times$ over
FlashAttention-2 and $1.23\times$ over SDPA.
The performance advantage is most pronounced on the ViT architecture, where \NAME{}
achieves up to $1.54\times$ speedup against FlexAttn.
In the FP8 precision, \NAME{} shows improvements with an average speedup of $1.16\times$
over FlexAttn (FP8) and $1.20\times$ over Tawa (FP8).

On the RTX 5090 platform (Table~\ref{tab:rq1-summary-5090}), \NAME{} continues achieving
higher relative throughput.
In the FP16 setting, \NAME{} achieves an average relative throughput of $1.26\times$ over
FlashAttention-2 and $1.42\times$ over SDPA.
% Similar to the GH200 results, the gains are substantial on the ViT architecture, where
% \NAME{} outperforms FlexAttention by $1.85\times$ and SDPA by $1.83\times$.
Compared to the TileLang baseline, \NAME{} shows a slight improvement, since \NAME{}
automatically selects TileLang as the optimal backend on the RTX 5090 architecture, and
uses the auto-scheduler and auto-tuner to improve kernel performance \mbox{beyond the
default configurations.}

\begin{table}[t!]
  \centering
  \small
  \setlength{\tabcolsep}{6.5pt}
  \caption{Geometric mean of relative throughput of \NAME{} against baselines on RTX 5090
  across sequence lengths.}
  \label{tab:rq1-summary-5090}
  \def\dgbox{\diagbox[width=6em, height=2.5em]{Model}{Method}}
  \renewcommand{\arraystretch}{0.93}
\begin{tabular}{l|cccc}
  \toprule
  \dgbox & Flash-2 & Flex & SDPA & TileLang \\
  \midrule
  GLM     & 1.16 & 1.10 & 1.74 & 1.01 \\
  LLAMA   & 1.32 & 1.14 & 1.59 & 1.01 \\
  QWEN2   & 1.28 & 1.14 & 1.33 & 1.00 \\
  QWEN3   & 1.32 & 1.15 & 1.33 & 1.00 \\
  VIT     & 1.23 & 1.30 & 1.18 & 1.02 \\
  \midrule
  \textbf{Average} & 1.26 & 1.16 & 1.42 & 1.01 \\
  \bottomrule
\end{tabular}
\end{table}

Figure~\ref{fig:gh200-throughput} and Figure~\ref{fig:5090-throughput} report the
detailed throughput for these models across varying sequence lengths
($1k$ to $32k$) and batch sizes, where the y-axis shows the throughput in TFLOPs per second.
\NAME{} consistently outperforms baselines because its auto-scheduler and auto-tuner
identify optimized parameters for every specific workload, whereas baselines rely on
static rules/configs.

Similarly, \NAME{} achieves greater throughput gains on ViT than on LLMs.
This is because ViT involves less common attention configurations, with fewer
layers, fewer attention heads, and smaller hidden dimensions, which the
baselines often do not optimize well. \NAME{}, instead, generates efficient
kernel optimizations for these workloads.
%\matei{Grammatically, the end of this paragraph sounds a bit weird. Normally, I'd
% suggest something but I am too tired so I will recommend asking someone who has a bit
% more sleep.} \yuchen{Further modified}

\begin{figure*}[t!]
  \centering
  % First Row: GH200 FP16
  \begin{subfigure}[b]{0.95\textwidth}
    \includegraphics[width=\textwidth]{figures/gh200-fp16-tc1.pdf}
    \vspace{-1.7em}
    \caption{GH200: FP16 performance (higher is better).}
    \label{fig:gh200-fp16-tc}
  \end{subfigure}

  % Second Row: GH200 FP8
  \begin{subfigure}[b]{0.95\textwidth}
    \includegraphics[width=\textwidth]{figures/gh200-fp8-tc1.pdf}
    \vspace{-1.7em}
    \caption{GH200: FP8 performance (higher is better).}
    \label{fig:gh200-fp8-tc}
  \end{subfigure}

  % Third Row: RTX 5090 FP16 (New)
  \begin{subfigure}[b]{0.95\textwidth}
    \includegraphics[width=\textwidth]{figures/rtx5090-fp16-tc1.pdf}
    \vspace{-1.7em}
    \caption{RTX 5090: FP16 performance (higher is better).}
    \label{fig:rtx5090-fp16-tc}
  \end{subfigure}

  \caption{Performance of \NAME{} kernels vs. kernels by other tensor compilers on GH200
    and RTX 5090 platforms.
  The y-axis shows relative performance of all compilers normalized to \NAME{} for each setup.}
  \vspace{-.1in}
  \label{fig:compiler-prefill-ablation}
\end{figure*}



\subsection{Detailed Kernel-level Latency Evaluation}

Figure~\ref{fig:compiler-prefill-ablation} illustrates the relative performance of
\NAME{} against leading tensor compilers across three attention-based operators on GH200
and RTX 5090 platforms.
Each subfigure reports the performance of baselines normalized to \NAME{} (y-axis) over
varying input sequence lengths (x-axis).
On the GH200 platform, \NAME{} achieves the best performance in the majority of setups.
At short sequence lengths, our auto-scheduler and auto-tuner gain a significant advantage
over Tawa by automatically selecting the faster TileLang~\cite{tilelang} backend and
tuning parameters towards smaller block sizes and fewer pipeline stages.
%
For instance, on short sequences, \NAME{} selects a lighter $64 \times 128$ tile size
with 4 warps, while Tawa defaults to a heavy $128 \times 128$ size with 8 warps.
This adaptability provides even larger benefits for FP8 precision. \NAME{} adjusts
parameters such as pipeline stages from 1 to 4 based on sequence length which the
performance of models in FP8 relies heavily on, so \NAME{} can achieve bigger performance
improvements compared to FP16.
%
On the RTX 5090, \NAME{} consistently outperforms baselines, with larger speedups
observed on longer sequences where our autoscheduler selects the rolling update
optimization to overcome TileLang.
In FP8 precision setups (Figure \ref{fig:gh200-fp8-tc}), \NAME{} demonstrates a
significant advantage, achieving \mbox{approximately} 20\% improvement on average over the
strongest baselines, enabled by the optimized \mbox{configurations} identified by our auto-tuner.
%
Our auto-tuner shows hardware adaptability by adjusting the search space parameters, e.g., block sizes, to accommodate the limited shared \mbox{memory of the RTX 5090.}

We do not show Neptune results since Neptune's performance is limited by Triton's performance,
due to being tied to Triton for code generation, which struggles to generate efficient code for these
operators on Hopper and Blackwell.
Our measurements show Neptune's speed is 0.5x-0.8x of that of Tawa or TileLang. %Table~\ref{tab:neptune_perf}
% show that Neptune is slower than other baselines.
This result indicates that the main bottleneck is backend code generation rather than schedule quality. Our results suggest that the schedules found by \NAME{} are already competitive with the manually written schedules used by the Neptune developers, and that the remaining gap comes from Triton's lower code quality on these platforms.
% \sasa{Separate: Somewhere need to say that our schedules are as good as those the Neptune developers manually wrote. }

% \sasa{What can we say here about specific cases? }

% \sasa{How can we relate the speedup to the success of our transformations/scheduling? }

% \sasa{We should also say why Neptune is not applicable here!! (Slow triton!)}

%\sasa{We should include some details about the schedules, which Matei mentioned in the
% slack "GQA and prefill attention are different in that GQA has an extra loop and some
% reshaping (which is merged into a block present in prefill-generated schedules). Other
% than that, only compiler flags/annotations differ at my level." Ofc emphasize the diffs
% and variety across schedules...}
% \sasa{Mention if there is even a small difference between schedules for GPU generations.}

\subsection{\NAME{} vs Manually-Written Kernels}

% Figure~\ref{fig:libs-prefill-comparison} compares \NAME{} against heavily optimized,
% hand-written kernel libraries on GH200 and RTX 5090 platforms.
% Across both hardware generations, \NAME{} generates kernels that consistently match or
% exceed the performance of state-of-the-art vendor libraries.
% %On GH200,
% \NAME{} matches the performance of the highly optimized cuDNN kernel.
% While hand-written libraries like cuDNN can be slightly faster in specific cases, they
% require intensive manual tuning that is difficult to maintain across different operators
% and hardware generations.

Appendix~\ref{sec:vendor-libs} compares \NAME{} against heavily optimized,
hand-written vendor libraries on GH200 and RTX 5090.
Across both hardware generations, \NAME{} matches or exceeds the performance of
state-of-the-art vendor libraries in most setups, and remains close to the highly
optimized cuDNN kernels in the remaining cases.

% For instance, supporting a new feature like SoftCap or ALiBi bias in cuDNN often
% requires waiting for vendor updates or accepting performance loss from unfused operations.
% In contrast, \NAME{} can automatically generate optimized, fused kernels for such custom variants.
%On the RTX 5090 (Blackwell) architecture, \NAME{} maintains consistent high performance
% across the entire range of configurations, matching or outperforming the
% \mbox{strongest hand-tuned baselines.}

% \sasa{We need a clear case here and explanation that even when cudnn is sometimes
% faster, it is pain to maintain and hand-tune...}

\subsection{Numerical Stability}
In Appendix~\ref{sec:num-stab}, we evaluate the numerical behavior of \NAME{}-generated kernels.
Using realistic Qwen2.5 inputs and an FP64 attention implementation as reference, \NAME{} shows numerical error comparable to Tri-Dao Attention and FlexAttention across PF Global, PF Causal, and PF GQA.
Averaged across the three operators, \NAME{} achieves an RMS absolute error of $4.96\times 10^{-5}$, compared with $4.90\times 10^{-5}$ for Tri-Dao and $5.02\times 10^{-5}$ for FlexAttention.
% \sasa{SAY WHICH Appendix. Also mention the averages in text. }



\subsection{Ablation Studies}

\begin{table}[t]
  \centering
  \caption{Ablation study of \NAME{}'s latency}
  \label{tab:ablation}
  \begin{tabular}{lc}
    \toprule
    \textbf{Configuration} & \textbf{Latency ($\mu s$)} \\
    \midrule
    \NAME{} (Full System) & 7.43 \\
    w/o Auto-tuning & 9.09 \\
    w/o Auto-tuning \& Expression Rewriting & 9.45 \\
    w/o Auto-scheduling & 10.93 \\
    \bottomrule
  \end{tabular}
\end{table}

To understand the importance of each component of \NAME{},
we perform ablation studies on the performance of \NAME{}'s kernels.
We select the Global attention operator with sequence length 256 on GH200.
Table \ref{tab:ablation} shows the results of this experiment.
The full \NAME{} system achieves the latency of 7.43 $\mu s$.
Disabling components progressively increases latency: removing the auto-tuner causes a
significant drop to 9.09$\mu$s. This shows the importance of parameter tuning.
Additionally disabling expression rewriting leads to further degradation to 9.45 $\mu s$.
Finally, the configuration with no auto-scheduling yields the highest latency of 10.93 $\mu s$.

\noindent\textbf{Autoscheduler Statistics.} \NAME{}'s auto-scheduling search is
relatively fast and finishes within a minute in most cases due to aggressive pruning. 
The auto-tuner takes a bit longer, taking from seconds to 10
minutes, as it needs to compile and runs many kernels.
\NAME-generated schedules contain a variety of different optimizations. GQA has several
reshape and broadcast operations added with respect to the base prefill attention, in
addition to an \mbox{extra loop nest.}
}

\section{\NA{Related Work}}
\label{sec:rw}
{\textbf{Tile-based Programming Frameworks.}
Tile languages such as Triton~\cite{triton}, TileLang~\cite{tilelang},
cuTile~\cite{cutile}, and others~\cite{cutlass,thunderkittens,tawa}
map tile computation to warps and threads but leave high-level optimizations,
such as fusion and expression rewriting, to programmers;
\NAME{} automates these and uses multiple tile compilers as backends.

\textbf{Scheduling Tensor Compilers and Auto-Schedulers.}
Schedule-based compilers such as Halide~\cite{halide}, TVM~\cite{tvm}, AKG~\cite{akg},
and Tiramisu~\cite{tiramisu} separate algorithms from schedules,
which auto-schedulers such as Ansor~\cite{ansor}, MetaSchedule~\cite{metaschedule},
and FlexTensor~\cite{flextensor} generate automatically.
Neptune~\cite{neptune} delegates low-level optimization to a tile compiler
and introduces advanced reduction fusion, but relies on manually written schedules;
\NAME{}'s auto-scheduler generates them automatically,
handling non-local advanced fusion while keeping programs tile-ready.

\textbf{Super-optimizing Tensor Compilers.}
Super-optimizers such as TASO~\cite{taso} and Mirage~\cite{mirage}
search over transformations that may not preserve semantics and verify candidates afterwards.
Mirage's multi-level search can take minutes to hours,
and its probabilistic correctness test can admit unsound kernels;
\NAME{} applies only semantics-preserving transformations
(up to floating-point reassociation).

\textbf{PyTorch-Based FlashAttention.}
Flashlight~\cite{flashlight} adds a semantic fusion pass to PyTorch's TorchInductor
to generate FlashAttention-style kernels from PyTorch attention definitions.
Its optimizations (semantic fusion and limited tiling) are graph rewrites
specific to TorchInductor, and it is tied to Triton;
\NAME{} combines operation-level, expression, and tile optimizations
in a three-IR pipeline over multiple tile compilers.

\textbf{Algebraic Expression Rewriting.}
Expression rewriting exploits algebraic identities, typically over whole tensors at the graph
level, as in TASO, TenSat~\cite{tensat}, and XLA~\cite{xla},
or at lower levels in TVM, MLIR~\cite{mlir}, and Lift~\cite{lift}.
\NAME{}'s VR-tile IR enables rewriting over tiles and across loop iterations,
and could adopt general-purpose rewriting engines such as egg~\cite{egg}.

\textbf{Graph-level Optimizers and LLM Serving.}
Graph compilers such as XLA, Glow~\cite{glow}, and BladeDISC~\cite{bladedisc}
fuse and transform whole computation graphs,
and LLM serving systems such as vLLM~\cite{vllm}, DeepSpeed Inference~\cite{deepspeed-inference},
FasterTransformer~\cite{fastertransformer}, and TensorRT-LLM~\cite{tensorrtllm}
manage KV caches and requests;
both rely on separate kernel generators or libraries, which \NAME{} can provide.

% \yifan{Update RW -- what came out in the last 3 months?}
}

\section{\NA{Conclusion}}
\label{sec:conclusion}
{We presented \NAME{}, a tensor compiler whose auto-scheduler automates advanced reduction fusion
on top of tile compilers, generating attention kernels from math-like specifications
that outperform state-of-the-art compilers and are on par with hand-optimized libraries.
}

\bibliographystyle{ACM-Reference-Format}
\bibliography{references}

\clearpage
\onecolumn
\appendix
\label{app}
\clearpage
\section*{Appendix}
\vspace{.3in}

\section{Experimental Details}
\label{sec:exp-details}

Table~\ref{tab:operators} lists the attention operators and the models they are extracted from,
and Table~\ref{tab:baselines} lists the versions and configurations of all baselines.
Each baseline supports a subset of the setups we evaluate.

For throughput, following~\cite{Megatron-LM}, we calculate the theoretical
floating-point operations (FLOPs) required for the attention computation
based on the model configuration (such as sequence length, number of heads, and hidden
dimension) and divide it by the measured latency.
We measure end-to-end latency with the CUDA Events API,
and kernel latency in operator experiments with Nsight Systems (\texttt{nsys})
for better precision on short-duration kernels.

\begin{table}[h]
  \small
  \centering
  \caption{LLM operators used in our evaluation. }
  \label{tab:operators}
  \begin{tabular}{l|l|l} \Xhline{2\arrayrulewidth}
    \textbf{Operator} & \textbf{Description}    & \textbf{Base Arch.} \\ \hline
    Global      & Global Attention        & ViT 1.2B \\
    Causal      & Causal Attention (MHA)  & \makecell[l]{Llama 2 7B, Qwen2 7B, \\ Qwen3 8B} \\
    GQA         & Grouped Query Attention & GLM-4 9B \\
    % \NA{Sparse} & \NA{Deepseek Sparse Attention} & \NA{DeepSeek-V3.2} \\
    \Xhline{2\arrayrulewidth}
  \end{tabular}
\end{table}

\begin{table}[h]
  \small
  \centering
  \caption{Baselines used in our evaluation.}
  \label{tab:baselines}
  \begin{tabular}{llll}
    \toprule
    \textbf{Baseline} & \textbf{Version} & \textbf{Kind} & \textbf{Configuration} \\
    \midrule
    Tawa~\cite{tawa}                    & 3.3            & Tile compiler & Stock attention, FlashAttention-3 strategy \\
    TileLang~\cite{tilelang}            & 0.1.8          & Tile compiler & Stock attention, FlashAttention-3 strategy \\
    OpenAI Triton~\cite{openai-fused-attn} & Triton 3.6.0 & Tile compiler & OpenAI's fused attention \\
    Tri-Dao Triton~\cite{flash-attn-repo}  & Triton 3.6.0 & Tile compiler & Dao-AILab's Triton attention \\
    Helion                              & 0.3.1          & Compiler      & \\
    FlexAttn                            & PyTorch 2.11.0 & Compiler      & \texttt{compile(mode=`max-autotune', backend=`inductor')} \\
    TVM~\cite{tvm}                      & 0.18           & Compiler      & \\
    \midrule
    cuDNN                               & 9.11.0         & Library       & \\
    Tri-Dao~\cite{flash-attn-repo}      & 2.8.3          & Library       & Dao-AILab's CUTLASS implementation \\
    FlashInfer~\cite{flashinfer}        & 0.6.7          & Library       & \\
    ThunderKittens~\cite{thunderkittens} & --            & Library       & Obtained from the ThunderKittens authors \\
    PyTorch~\cite{pytorch}              & 2.11.0         & Library       & \texttt{compile(mode=`max-autotune', backend=`inductor')} \\
    \bottomrule
  \end{tabular}
\end{table}

\clearpage

\section{\NAME{} vs Vendor Libraries}
\label{sec:vendor-libs}

Appendix~\ref{sec:vendor-libs} provides the full comparison of Nautilus against
heavily optimized vendor libraries on GH200 and RTX 5090,
including detailed results for different attention variants and
sequence lengths.



\begin{figure*}[hb!]
  \centering
  % First Row: GH200 Baseline
  \begin{subfigure}[b]{0.95\textwidth}
    \includegraphics[width=\textwidth]{figures/gh200-fp16-baseline1.pdf}
    \vspace{-1.7em}
    \caption{GH200: FP16 performance (higher is better).}
    \label{fig:gh200-fp16-libs}
  \end{subfigure}

  % Second Row: RTX 5090 Baseline (New)
  \begin{subfigure}[b]{0.95\textwidth}
    \includegraphics[width=\textwidth]{figures/rtx5090-fp16-baseline1.pdf}
    \vspace{-1.7em}
    \caption{RTX 5090: FP16 performance (higher is better).}
    \label{fig:rtx5090-fp16-libs}
  \end{subfigure}

  \caption{Performance of \NAME{} kernels vs. vendor libraries on GH200 and RTX 5090 platforms.
  The y-axis shows relative performance normalized to \NAME{} for each setup.}
  \label{fig:libs-prefill-comparison}
  %\vspace{-.15in}
\end{figure*}

\clearpage



\section{Numerical Stability of \NAME{} Kernels}
\label{sec:num-stab}

\noindent
\begin{minipage}[t]{0.46\textwidth}
\vspace{0pt}
Rewriting or re-associating floating-point operations can affect a kernel's numerical
behavior. To study this effect, we compare \NAME{}, Tri-Dao, and FlexAttention on
different attention operators. The results show that \NAME{} has competitive numerical
behavior. For all three systems, we use a common set of inputs and compare their outputs
to an FP64 unfused attention kernel baseline.
We keep the schedule and configuration of each kernel the same as in the performance
evaluation. To collect realistic attention inputs, we run a pretrained Qwen2.5 model
on 1000 WikiText sentences with sequence length 2048, and record the resulting input tensors.

In Table~\ref{tab:numeric_error_absolute_error}, we report the numerical error of the
three frameworks using root mean square error (RMS) together with the 90th and 99th
percentile error magnitudes. We also examine how per-sample (per-sentence) RMSE is
distributed. The results for PF Global, PF Causal, and PF GQA are shown in
Figure~\ref{fig:rmse-kde-all}. Overall, Nautilus stays very close to Tri-Dao
Attention in error and is consistently better than FlexAttention.
\end{minipage}%
\hfill
\begin{minipage}[t]{0.5\textwidth}
\vspace{0pt}
\centering\footnotesize\setlength{\tabcolsep}{4pt}

\captionof{table}{Numerical Error Summary for FP16 (Absolute Error)}
\label{tab:numeric_error_absolute_error}
\begin{tabular}{llrrr}
\hline
Method & Variant & RMS & 90th \% & 99th \% \\
\hline
\multirow{3}{*}{Nautilus}
& PF Global & 4.9e-05 & 5.1e-05 & 5.5e-05 \\
& PF Causal & 5.0e-05 & 5.1e-05 & 5.3e-05 \\
& PF GQA    & 5.0e-05 & 5.1e-05 & 5.3e-05 \\
\hline
\multirow{3}{*}{Tri-Dao Attention}
& PF Global & 4.8e-05 & 5.0e-05 & 5.1e-05 \\
& PF Causal & 4.9e-05 & 5.0e-05 & 5.2e-05 \\
& PF GQA    & 4.9e-05 & 5.0e-05 & 5.2e-05 \\
\hline
\multirow{3}{*}{Flex Attention}
& PF Global & 5.0e-05 & 5.2e-05 & 5.4e-05 \\
& PF Causal & 5.0e-05 & 5.1e-05 & 5.3e-05 \\
& PF GQA    & 5.0e-05 & 5.1e-05 & 5.3e-05 \\
\hline
\end{tabular}

\vspace{1em}
\captionof{table}{Numerical Error Summary for BF16 (Absolute Error)}
\label{tab:numeric_error_absolute_error_bf16}
\begin{tabular}{llrrr}
\hline
Method & Variant & RMS & 90th \% & 99th \% \\
\hline
\multirow{3}{*}{Nautilus}
& PF Global & 4.42e-04 & 4.45e-04 & 4.49e-04 \\
& PF Causal & 3.84e-04 & 3.87e-04 & 3.91e-04 \\
& PF GQA    & 3.83e-04 & 3.87e-04 & 3.90e-04 \\
\hline
\multirow{3}{*}{Tri-Dao Attention}
& PF Global & 4.41e-04 & 4.45e-04 & 4.48e-04 \\
& PF Causal & 3.84e-04 & 3.87e-04 & 3.91e-04 \\
& PF GQA    & 3.84e-04 & 3.87e-04 & 3.91e-04 \\
\hline
\multirow{3}{*}{Flex Attention}
& PF Global & 4.42e-04 & 4.45e-04 & 4.49e-04 \\
& PF Causal & 3.84e-04 & 3.87e-04 & 3.91e-04 \\
& PF GQA    & 3.84e-04 & 3.87e-04 & 3.91e-04 \\
\hline
\multirow{3}{*}{Flashlight}
& PF Global & 4.42e-04 & 4.45e-04 & 4.49e-04 \\
& PF Causal & 3.84e-04 & 3.87e-04 & 3.91e-04 \\
& PF GQA    & 3.84e-04 & 3.87e-04 & 3.91e-04 \\
\hline
\end{tabular}

\vspace{1em}
\captionof{table}{Numerical Error Summary for FP8 (Absolute Error)}
\label{tab:numeric_error_absolute_error_fp8}
\begin{tabular}{llrrr}
\hline
Method & Variant & RMS & 90th \% & 99th \% \\
\hline
\multirow{3}{*}{Nautilus}
& PF Global & 1.11e-02 & 1.28e-02 & 1.46e-02 \\
& PF Causal & 9.89e-03 & 1.12e-02 & 1.34e-02 \\
& PF GQA    & 9.88e-03 & 1.12e-02 & 1.34e-02 \\
\hline
\multirow{3}{*}{Tri-Dao Attention}
& PF Global & 1.16e-02 & 1.33e-02 & 1.50e-02 \\
& PF Causal & 1.02e-02 & 1.15e-02 & 1.37e-02 \\
& PF GQA    & 1.03e-02 & 1.16e-02 & 1.38e-02 \\
\hline
\multirow{3}{*}{Flex Attention}
& PF Global & 1.55e-02 & 1.68e-02 & 1.85e-02 \\
& PF Causal & 1.39e-02 & 1.50e-02 & 1.67e-02 \\
& PF GQA    & 1.39e-02 & 1.50e-02 & 1.67e-02 \\
\hline
\end{tabular}
\end{minipage}

\begin{figure}[h]
  \centering
\vspace{.2in}
  \begin{subfigure}[t]{0.32\textwidth}
    \centering
    \includegraphics[width=\linewidth]{figures/qwen_rmse_kde_3curves_1000_global.pdf}
    \caption{PF Global.}
    \label{fig:rmse-kde-global}
  \end{subfigure}
  \hfill
  \begin{subfigure}[t]{0.32\textwidth}
    \centering
    \includegraphics[width=\linewidth]{figures/qwen_rmse_kde_3curves_1000_causal.pdf}
    \caption{PF Causal.}
    \label{fig:rmse-kde-causal}
  \end{subfigure}
  \hfill
  \begin{subfigure}[t]{0.32\textwidth}
    \centering
    \includegraphics[width=\linewidth]{figures/qwen_rmse_kde_3curves_1000_gqa.pdf}
    \caption{PF GQA.}
    \label{fig:rmse-kde-gqa}
  \end{subfigure}

  \caption{Kernel numerical stability comparison on Qwen across three attention settings for FP16.}
  \label{fig:rmse-kde-all}
\end{figure}

\begin{figure}[h]
  \centering
\vspace{.2in}
  \begin{subfigure}[t]{0.32\textwidth}
    \centering
    \includegraphics[width=\linewidth]{figures/bf16_prefill_global_rmse.pdf}
    \caption{PF Global.}
    \label{fig:rmse-kde-global}
  \end{subfigure}
  \hfill
  \begin{subfigure}[t]{0.32\textwidth}
    \centering
    \includegraphics[width=\linewidth]{figures/bf16_prefill_causal_rmse.pdf}
    \caption{PF Causal.}
    \label{fig:rmse-kde-causal}
  \end{subfigure}
  \hfill
  \begin{subfigure}[t]{0.32\textwidth}
    \centering
    \includegraphics[width=\linewidth]{figures/bf16_prefill_gqa_rmse.pdf}
    \caption{PF GQA.}
    \label{fig:rmse-kde-gqa}
  \end{subfigure}

  \caption{Kernel numerical stability comparison on Qwen across three attention settings for BF16.}
  \label{fig:rmse-kde-all_bf16}
\end{figure}

\begin{figure}[h]
  \centering
\vspace{.2in}
  \begin{subfigure}[t]{0.32\textwidth}
    \centering
    \includegraphics[width=\linewidth]{figures/fp8_prefill_global_rmse.pdf}
    \caption{PF Global.}
    \label{fig:rmse-kde-global}
  \end{subfigure}
  \hfill
  \begin{subfigure}[t]{0.32\textwidth}
    \centering
    \includegraphics[width=\linewidth]{figures/fp8_prefill_causal_rmse.pdf}
    \caption{PF Causal.}
    \label{fig:rmse-kde-causal}
  \end{subfigure}
  \hfill
  \begin{subfigure}[t]{0.32\textwidth}
    \centering
    \includegraphics[width=\linewidth]{figures/fp8_prefill_gqa_rmse.pdf}
    \caption{PF GQA.}
    \label{fig:rmse-kde-gqa}
  \end{subfigure}

  \caption{Kernel numerical stability comparison on Qwen across three attention settings for FP8.}
  \label{fig:rmse-kde-all_bf16}
\end{figure}


\begin{table*}[t]
\centering
\small
\caption{Scheduling primitives used by the Nautilus auto-scheduler.}
\label{tab:primitives}
\begin{tabular}{clll}
\hline
Abbr. & Primitive & Short Description & Source \\
\hline
TL  & TileLoop($l_i$, $s_i$) & Split each $l_i$ into two loops with split factor $s_i$, and reorder these loops & TVM \\
BBI & BindBlockIdx($l$, \texttt{blockIdx}) & Fuse loops $l$ and associate them with a GPU thread block dimension & TVM \\
\hline
CRW & CacheReadWrite($n$, $t$, \texttt{place}) & Copy regions of $t$ accessed by $n$ from its current placement to \texttt{place} & TVM \\
SP  & SetPlacement($t$, \texttt{place}) & Set the placement of $t$ to \texttt{place}, affecting all accesses to $t$ & TVM \\
\hline
CI  & ComputeInline($n$, \texttt{dir}) & \begin{tabular}[c]{@{}l@{}}Inline the computation in $n$ into all producers/consumers,\\ and remove $n$ from the block graph\end{tabular} & TVM \\
CA  & ComputeAt($n$, $l$, \texttt{dir}) & Loop transformation-based operator fusion: fuse $n$ under $l$ & TVM \\
RU  & RollingUpdate($n_r$, $l$) & Advanced reduction operator fusion: fuse reduction $n_r$ under $l$ & Neptune \\
SKU & SplitKUpdate($n_r$, $l$) & \begin{tabular}[c]{@{}l@{}}Split reduction $n_r$ into local and global reductions,\\ and fuse the local reduction under $l$\end{tabular} & Neptune \\
\hline
DR  & DecomposeReduction($n_r$, $l$) & Move the initialization of reduction $n_r$ outside of $l$ & TVM \\
SSB & SplitScanBuffer($n_s$, $l$) & \begin{tabular}[c]{@{}l@{}}Multi-buffer the output of scan operation $n_s$\\ and update the buffers per $l$ iteration\end{tabular} & Neptune \\
\hline
PA  & PadAxis($n$, $t$, $a$, $e$) & \begin{tabular}[c]{@{}l@{}}Pad axis $a$ of $t$ to extent $e$ through every statement that carries it,\\ zero-filling inputs and cropping outputs\end{tabular} & Nautilus \\
PIE & PropagateIfThenElse($n$, $l$) & Narrow $l$ to the iterations where the mask in $n$ can hold & Neptune \\
FA  & FoldAccumulators($n$) & \begin{tabular}[c]{@{}l@{}}Fold terms added to a matmul in $n$ (bias, mask, other matmuls)\\ into its accumulator's initial value\end{tabular} & Nautilus \\
RB  & ReorderBlocks($l$, $o$) & Reorder the statements in the body of $l$ to legal order $o$ & Nautilus \\
\hline
\end{tabular}
\end{table*}



\begin{table*}[t]
\centering
\small
\caption{Tuning decisions in Nautilus schedules beyond the primitives' arguments. Structural choices fork
schedule seeds; sampled ones are recorded in the trace, so the auto-tuner mutates them and children
inherit good values.}
\label{tab:decisions}
\begin{tabular}{llll}
\hline
Decision & Values & Kind & Constraint / Note \\
\hline
Fusion strategy per reduction & classic, RU, SKU, RU+SKU hybrid & structural & \begin{tabular}[c]{@{}l@{}}Hybrid (3-level split) needs $\geq$16 rows in the tile\end{tabular} \\
Row padding                   & padded to 16, unpadded          & structural & Only for matmuls with $<$16 query rows (decode, MTP) \\
Tile loop sizes               & $s_i \in [32, 256]$             & sampled    & \begin{tabular}[c]{@{}l@{}}Accumulator side $\leq$128 and $M\!\times\!N \leq$ 1/4 of the register file;\\ reduction left whole if its operands fit half of shared memory\end{tabular} \\
Split count (hybrid)          & $[2,\ \mathrm{KV\ tiles}/2]$       & sampled    & Number of thread blocks splitting the KV axis \\
Warps per block               & 2, 4, 8                         & sampled    & Shared by all backends \\
Pipeline stages               & 0, 1, 2 (TileLang); 1, 2, 3 (Triton/Tawa) & sampled & \begin{tabular}[c]{@{}l@{}}On the rolling-update loop; lowered until the\\ multi-buffered tiles fit in shared memory\end{tabular} \\
Warp specialization           & off, on                         & sampled    & TileLang, sm\_90 and later \\
TMA copies                    & off, on                         & sampled    & TileLang, sm\_90 and later \\
Accumulator folding (FA)      & off, on                         & sampled    & Faster at small warp counts only \\
Statement order (RB)          & any legal order                 & sampled    & Uniform over orders keeping RAW/WAR/WAW dependencies \\
\hline
\end{tabular}
\end{table*}

\end{document}
